\section{Experiments}
\label{sec:experiments}

We evaluate the self-evolving framework along five axes: whether it consistently improves both understanding and generation across architecturally distinct models (Sec.~\ref{sec:main_results}), how it compares to existing self-supervised post-training methods under equivalent supervision constraints (Sec.~\ref{sec:comparison}), which components are essential to its effectiveness (Sec.~\ref{sec:ablation}), whether the training dynamics exhibit the expected curriculum escalation and capability progression (Sec.~\ref{sec:dynamics}), and how the improvements manifest qualitatively in generated content and question complexity (Sec.~\ref{sec:qualitative}).

\subsection{Experimental Setup}
\label{sec:setup}

% TODO: Fill with actual experimental details once training is complete.

\noindent\textbf{Models.}
We apply the framework to three unified understanding and generation models that employ fundamentally different image synthesis mechanisms: BLIP3o-8B (diffusion-based denoising), BAGEL (flow-matching with a variational autoencoder), and VARGPT (autoregressive discrete token generation). For each model, we attach three LoRA adapters (proposer, solver, generator) to the frozen backbone with rank $r$ and train them using identical hyperparameters and training schedules.

\noindent\textbf{Training data.}
All training uses unlabeled images only. No captions, question-answer pairs, or human annotations are used at any stage. The image sources and quantities are reported per model.

\noindent\textbf{Understanding benchmarks.}
We evaluate visual understanding on MMMU, MMBench, TextVQA, SEED-Bench, RealWorldQA (RWQA), MM-Vet, and both perception (MME-P) and cognition (MME-C) splits of MME.

\noindent\textbf{Generation benchmarks.}
We evaluate image generation on GenEval (compositional attribute accuracy), DPG-Bench (dense prompt following), and WISE (open-ended generation quality).

\noindent\textbf{Training details.}
% TODO: Fill with learning rate, batch size, number of steps, LoRA rank,
% number of prompt perturbations N, STE window size, U:G step ratio,
% and other hyperparameters.


Table~\ref{tab:qualitative_comparison} provides a structured qualitative comparison of our framework against existing post-training methods along the five design axes introduced in Section~\ref{sec:related_work}. Our approach is the only method that simultaneously requires no external supervision, optimizes both understanding and generation bidirectionally, employs multiple complementary reward signals, incorporates explicit stability controls, and demonstrates generality across three architecturally distinct model families.

% =====================================================================
% Table 1: Structured Qualitative Comparison (5 Axes from Related Work)
% =====================================================================
\begin{table}[tb]
  \caption{Structured comparison of self-evolving and post-training approaches for unified understanding and generation models along five design axes. ``Ext. Sup.'' denotes reliance on external supervision such as curated labels, SFT data, or reward models. ``Bidir.'' denotes whether both understanding and generation are jointly optimized. ``Reward'' indicates the nature of the training signal: single scalar or multi-component. ``Stability'' denotes the presence of explicit collapse prevention mechanisms. ``Arch. Dep.'' denotes whether the method is tied to a specific model architecture. Entries are derived from each method's published description.}
  \label{tab:qualitative_comparison}
  \centering
  \small
  \resizebox{\linewidth}{!}{%
  \begin{tabular}{@{}lccccc@{}}
    \toprule
    Method & Ext. Sup. & Bidir. & Reward & Stability & Arch. Dep. \\
    \midrule
    UniCorn \cite{han2026unicorn}              & Judge training & Yes & Single (judge)           & Not reported      & Single family \\
    SEER \cite{tang2026endogenous}             & Reward activation (RLVR) & U only & Single (evaluator)  & Reprompter loop   & Single family \\
    UniGame \cite{su2025unigame}               & None           & Yes & Single (adversarial)     & Perturbation norm & Architecture-specific \\
    SRUM \cite{jin2025srum}                    & Reward design  & G from U & Multi (global-local)& Manual granularity& Single family \\
    EvoLMM \cite{thawakar2025evolmm}           & None           & U only & Single (self-consistency)& Curriculum        & Single family \\
    UAE \cite{yan2025understandinggeneration}   & Supervised init. & Yes & Single (reconstruction) & Multi-stage       & Encoder-decoder \\
    RecA \cite{xie2025reconstruction}          & None           & G from U & Single (reconstruction) & Not reported  & Architecture-specific \\
    Internal Gap \cite{han2025internalgap}     & SFT + DPO data & Yes & Single (gap score)       & Curriculum SFT    & Single family \\
    UniRL \cite{mao2025unirl}                  & SFT data       & Yes & Task-specific            & Iterative SFT     & Single family \\
    CoRL \cite{jiang2025coreinforcement}       & Task-specific rewards & Yes & Task-specific       & Multi-stage RL    & Single family \\
    SUDER \cite{hong2025dualselfrewards}       & None           & Yes & Single (dual likelihood) & Not reported      & Reversible interface \\
    SILMM \cite{qu2024silmm}                  & Self-feedback  & G only & Single (DPO pref.)      & Iterative         & Single family \\
    ILLUME \cite{wang2024illume}               & Self-assessment & Yes & Single (consistency)    & Not reported      & Single family \\
    UnifiedReward \cite{wang2025unifiedreward} & Reward model data & Enabler & Multi (cross-task) & Not reported      & Model-agnostic \\
    \midrule
    \textbf{Ours}                              & \textbf{None}  & \textbf{Yes} & \textbf{Multi (SC+STE+QA+Cyc)} & \textbf{EMA + buckets + filter} & \textbf{3 families} \\
    \bottomrule
  \end{tabular}%
  }
\end{table}


\subsection{Main Results}
\label{sec:main_results}

Table~\ref{tab:main_results} presents understanding and generation results across all three backbone architectures. Existing unified understanding and generation models (Janus-Pro-7B, SEED-X, Emu3) are included as reference points to contextualize the performance landscape. Each of our models is evaluated in its pretrained base configuration and after applying the self-evolving framework, with the absolute improvement ($\Delta$) reported for each benchmark.

% TODO: Add discussion of results once numbers are available.
% Key points to cover:
% - Consistent improvement on both understanding AND generation across all 3 models
% - How our post-trained models compare to other unified models
% - Highlight that the same algorithm/hyperparameters work across diffusion, flow-matching, and autoregressive
% - Note any benchmark where improvements are particularly large or small, and why
% - Generation Avg column provides a single summary statistic

% =====================================================================
% Table 2: Primary Comparison — Self-Supervised Methods (Fair Comparison)
% Only methods that require NO curated SFT data, NO external reward
% models, and NO supervised initialization are included here.
% =====================================================================
\begin{table}[tb]
  \caption{Comparison with self-supervised post-training methods using \emph{relative gains over each method's own base model}. All methods listed here operate without curated labels, supervised fine-tuning data, or external reward models, ensuring a fair comparison under equivalent supervision constraints. Understanding and generation benchmarks are separated by a vertical rule.}
  \label{tab:selfsup_comparison}
  \centering
  \small
  \resizebox{\linewidth}{!}{%
  \begin{tabular}{@{}p{2.4cm}p{2.0cm}cccccccc|ccc@{}}
    \toprule
    & & \multicolumn{8}{c|}{Understanding} & \multicolumn{3}{c}{Generation} \\
    \cmidrule(lr){3-10} \cmidrule(l){11-13}
    Method & Backbone & MMMU  & MMBench & TextVQA & SEED & RWQA & MMVet & MME-P & MME-C & GenEval & DPG-Bench & WISE \\
    \midrule
    UniGame \cite{su2025unigame} & UMM & -- & -- & -- & -- & -- & -- & -- & -- & -- & -- & -- \\
    SUDER \cite{hong2025dualselfrewards} & Janus-Pro-7B & -- & -- & -- & -- & -- & -- & -- & -- & -- & -- & -- \\
    EvoLMM \cite{thawakar2025evolmm} & Janus-Pro-7B & -- & -- & -- & -- & -- & -- & -- & -- & -- & -- & -- \\
    SILMM \cite{qu2024silmm} & SEED-X & -- & -- & -- & -- & -- & -- & -- & -- & -- & -- & -- \\
    ILLUME \cite{wang2024illume} & SEED-X & -- & -- & -- & -- & -- & -- & -- & -- & -- & -- & -- \\
    \textbf{Ours} & BLIP3o-8B & -- & -- & -- & -- & -- & -- & -- & -- & -- & -- & -- \\
    \bottomrule
  \end{tabular}%
  }
\end{table}


% =====================================================================
% Table 3: Extended Reference — Supervised Methods (NOT directly comparable)
% These methods use curated SFT data, external reward models, or
% supervised initialization. Shown as reference upper bounds.
% =====================================================================
\begin{table}[tb]
  \caption{Reference comparison with methods that use additional supervision. These approaches employ curated fine-tuning data, external reward models, or supervised initialization stages, and are therefore not directly comparable to our fully unsupervised framework. Results are shown as reference points to contextualize the performance gap between supervised and self-supervised post-training.}
  \label{tab:supervised_reference}
  \centering
  \small
  \resizebox{\linewidth}{!}{%
  \begin{tabular}{@{}p{2.4cm}p{2.0cm}p{3.0cm}cccccccc|ccc@{}}
    \toprule
    & & & \multicolumn{8}{c|}{Understanding} & \multicolumn{3}{c}{Generation} \\
    \cmidrule(lr){4-11} \cmidrule(l){12-14}
    Method & Backbone & Ext. Supervision & MMMU  & MMBench & TextVQA & SEED & RWQA & MMVet & MME-P & MME-C & GenEval & DPG-Bench & WISE \\
    \midrule
    UniRL \cite{mao2025unirl} & Janus & SFT data & -- & -- & -- & -- & -- & -- & -- & -- & -- & -- & -- \\
    UniCorn \cite{han2026unicorn} & UMM & Judge training & -- & -- & -- & -- & -- & -- & -- & -- & -- & -- & -- \\
    CoRL \cite{jiang2025coreinforcement} & Janus-Pro-1B & Task-specific rewards & -- & -- & -- & -- & -- & -- & -- & -- & -- & -- & -- \\
    UAE \cite{yan2025understandinggeneration} & Janus-Pro-7B & Supervised init. & -- & -- & -- & -- & -- & -- & -- & -- & -- & -- & -- \\
    Internal Gap \cite{han2025internalgap} & UMM & SFT + DPO data & -- & -- & -- & -- & -- & -- & -- & -- & -- & -- & -- \\
    \midrule
    \textbf{Ours} & BLIP3o-8B & \textbf{None} & -- & -- & -- & -- & -- & -- & -- & -- & -- & -- & -- \\
    \bottomrule
  \end{tabular}%
  }
\end{table}


% =====================================================================
% Table 4: Model-Agnostic Understanding Benchmarks
% =====================================================================
\begin{table}[tb]
  \caption{Understanding benchmarks across three architecturally distinct unified models. Each model is evaluated in its base configuration and after applying our self-evolving framework, with the difference ($\Delta$) reported.}
  \label{tab:model_agnostic_understanding}
  \centering
  \small
  \resizebox{\linewidth}{!}{%
  \begin{tabular}{@{}p{2.8cm}cccccccc@{}}
    \toprule
    Model & MMMU  & MMBench & TextVQA & SEED & RWQA & MMVet & MME-P & MME-C \\
    \midrule
    BLIP3o-8B (Base) & -- & -- & -- & -- & -- & -- & -- & -- \\
    BLIP3o-8B (Ours) & -- & -- & -- & -- & -- & -- & -- & -- \\
    $\Delta$  & -- & -- & -- & -- & -- & -- & -- & -- \\
    \midrule
    BAGEL (Base) & -- & -- & -- & -- & -- & -- & -- & -- \\
    BAGEL (Ours) & -- & -- & -- & -- & -- & -- & -- & -- \\
    $\Delta$ & -- & -- & -- & -- & -- & -- & -- & -- \\
    \midrule
    VARGPT (Base) & -- & -- & -- & -- & -- & -- & -- & -- \\
    VARGPT (Ours) & -- & -- & -- & -- & -- & -- & -- & -- \\
    $\Delta$ & -- & -- & -- & -- & -- & -- & -- & -- \\
    \bottomrule
  \end{tabular}%
  }
\end{table}


% =====================================================================
% Table 5: Model-Agnostic Generation Benchmarks
% =====================================================================
\begin{table}[tb]
  \caption{Generation benchmarks across three architecturally distinct unified models. Each model is evaluated in its base configuration and after applying our self-evolving framework, with the difference ($\Delta$) reported.}
  \label{tab:model_agnostic_generation}
  \centering
  \small{%
  \begin{tabular}{lcccc}
    \toprule
    Model & GenEval & DPG-Bench & WISE & Avg \\
    \midrule
    BLIP3o (Base) & -- & -- & -- & -- \\
    BLIP3o (Ours) & -- & -- & -- & -- \\
    $\Delta$ & -- & -- & -- & -- \\
    \midrule
    BAGEL (Base) & -- & -- & -- & -- \\
    BAGEL (Ours) & -- & -- & -- & -- \\
    $\Delta$ & -- & -- & -- & -- \\
    \midrule
    VARGPT (Base) & -- & -- & -- & -- \\
    VARGPT (Ours) & -- & -- & -- & -- \\
    $\Delta$ & -- & -- & -- & -- \\
    \bottomrule
  \end{tabular}%
  }
\end{table}


% =====================================================================
% Table 6: Ablation Study
% =====================================================================
\begin{table}[tb]
  \caption{Ablation study on BLIP3o-8B measuring the contribution of individual framework components. Each row removes or replaces one component while keeping all others at their default configuration. Understanding is measured by the average across MMMU, MMBench, SEED, and MMVet. Generation is measured by GenEval.}
  \label{tab:ablation}
  \centering
  \small
  \begin{tabular}{lcc}
    \toprule
    Configuration & Und. Avg & GenEval \\
    \midrule
    Full framework & -- & -- \\
    \midrule
    \multicolumn{3}{l}{\emph{Understanding signal ablations}} \\
    \quad w/o STE (self-consistency only) & -- & -- \\
    \quad w/o self-consistency (STE only) & -- & -- \\
    \quad w/o prompt perturbation (temperature sampling) & -- & -- \\
    \quad w/o adaptive entropy target & -- & -- \\
    \midrule
    \multicolumn{3}{l}{\emph{Generation signal ablations}} \\
    \quad w/o QA fidelity scoring & -- & -- \\
    \quad w/o cycle-consistent captioning & -- & -- \\
    \quad w/o diversity reward & -- & -- \\
    \midrule
    \multicolumn{3}{l}{\emph{Optimization ablations}} \\
    \quad Proposer: REINFORCE instead of GRPO & -- & -- \\
    \quad w/o bucket-stratified baselines & -- & -- \\
    \quad w/o quality certificate filter & -- & -- \\
    \midrule
    \multicolumn{3}{l}{\emph{Architecture ablations}} \\
    \quad w/o generation loop (understanding only) & -- & -- \\
    \quad w/o understanding loop (generation only) & -- & -- \\
    \bottomrule
  \end{tabular}
\end{table}



\subsection{Comparison with Self-Supervised Methods}
\label{sec:comparison}

Table~\ref{tab:selfsup_comparison} provides a comprehensive comparison organized into three groups. Base unified models, including the backbones used by competing methods and our own three architectures, establish the starting point before any post-training. Self-supervised post-training methods, which operate without curated labels, supervised fine-tuning data, or external reward models, constitute the fair comparison group. Our results are reported across all three backbones to demonstrate that the improvements are not specific to a single architecture. Methods that employ additional supervision are reported separately in the supplementary material as reference upper bounds.

% TODO: Add discussion once numbers are filled.
% Key points to cover:
% - How our method compares to each self-supervised baseline
% - Highlight that we achieve competitive or superior results without ANY external signal
% - Note backbone differences where relevant (our BLIP3o-8B vs their Janus-Pro-7B, etc.)
% - Our three base models serve as direct ablation: base vs ours shows the effect of our framework
% - Reference supplementary Table S1 for supervised comparison


\subsection{Ablation Study}
\label{sec:ablation}

We ablate each framework component on BLIP3o-8B to isolate its contribution. Table~\ref{tab:ablation} reports understanding performance (averaged across MMMU, MMBench, SEED, and MMVet) and generation performance (GenEval) for each configuration.

% TODO: Add discussion once numbers are filled.
% Key points to cover:
% - Understanding signal ablations: STE vs self-consistency, importance of prompt perturbation
% - Generation signal ablations: QA fidelity vs cycle consistency complementarity
% - Optimization ablations: GRPO vs REINFORCE, bucket-stratified baselines
% - Architecture ablations: joint vs understanding-only vs generation-only
%   (this directly tests whether the two loops benefit each other)


\subsection{Training Dynamics}
\label{sec:dynamics}

Figure~\ref{fig:dynamics} visualizes the evolution of key training signals over the course of optimization, providing direct evidence that the self-evolving loop operates as intended.

\begin{figure*}[tb]
  \centering
  % TODO: Replace with actual training dynamics plots.
  % Recommended layout: 2x2 grid of subplots.
  % \includegraphics[width=\linewidth]{figures/training_dynamics.pdf}
  \fbox{\parbox{0.95\linewidth}{\centering\vspace{3.5cm}\textbf{[Training Dynamics Figure Placeholder]}\\\smallskip
  Four subplots: (a) Mean STE difficulty over training steps, (b) Self-consistency entropy distribution,\\
  (c) Easy/Medium/Hard bucket proportions, (d) Solver validation accuracy.\vspace{3.5cm}}}
  \caption{Training dynamics of the self-evolving loop on BLIP3o-8B. (a) Mean STE difficulty increases steadily as the proposer learns to generate harder questions, confirming progressive curriculum escalation. (b) Self-consistency entropy shifts from a zero-dominated distribution toward higher values as the proposer targets the solver's competence boundary. (c) The proportion of questions in the easy bucket decreases over training while medium and hard buckets grow, demonstrating effective difficulty calibration. (d) Solver validation accuracy on held-out benchmarks improves monotonically, confirming that harder questions translate to genuine capability gains rather than spurious difficulty inflation.}
  \label{fig:dynamics}
\end{figure*}

% TODO: Add discussion once plots are generated.
% Key points to cover:
% - (a) STE difficulty increases: proposer generates harder questions over time
% - (b) Entropy distribution shifts: early training dominated by H_sc=0, later more spread
% - (c) Bucket proportions: easy decreases, medium/hard increase — curriculum works
% - (d) Solver accuracy: monotonic improvement validates that difficulty escalation is productive
% - Optionally mention generation reward trend if space allows


\subsection{Qualitative Analysis}
\label{sec:qualitative}

Figure~\ref{fig:qualitative} presents representative examples that illustrate the framework's impact on both understanding and generation quality.

\begin{figure*}[tb]
  \centering
  % TODO: Replace with actual qualitative comparison figure.
  % Recommended layout: Top row = generation (real image, base generated, ours generated, prompt),
  %                     Bottom row = proposer question evolution (early vs late training).
  % \includegraphics[width=\linewidth]{figures/qualitative_results.pdf}
  \fbox{\parbox{0.95\linewidth}{\centering\vspace{4cm}\textbf{[Qualitative Results Figure Placeholder]}\\\smallskip
  Top: Generation comparison (prompt $\rightarrow$ base model output vs.\ ours) across diverse scenes.\\
  Bottom: Proposer question evolution showing early-stage simple questions vs.\ late-stage complex questions.\vspace{4cm}}}
  \caption{Qualitative comparison of generation and understanding before and after self-evolving training. Top two rows: given a text prompt, we compare images generated by the base model and our self-evolved model across three architectures. The self-evolved models produce images with more accurate attribute binding, improved spatial coherence, and better prompt adherence. Bottom row: proposer question evolution from early to late training, showing progression from simple object identification questions to compositionally complex queries that probe spatial relationships, counting, and contextual reasoning.}
  \label{fig:qualitative}
\end{figure*}

% TODO: Add discussion once qualitative examples are collected.
% Key points to cover:
% - Generation: specific improvements in attribute accuracy (color, count, spatial layout)
% - Generation: holistic scene coherence improvements
% - Proposer evolution: concrete examples of question complexity increasing
% - Optionally: 1-2 failure cases showing current limitations
